\documentclass[11pt]{article}
\usepackage[margin=1in]{geometry}
\usepackage{amsmath,amssymb}
\usepackage{hyperref}
\hypersetup{colorlinks=true,linkcolor=blue,citecolor=blue,urlcolor=blue}

\newcommand{\rd}{\mathrm{d}}
\newcommand{\pd}[2]{\frac{\partial #1}{\partial #2}}
\newcommand{\lap}{\nabla^2}
\newcommand{\bv}{\mathbf{v}}
\newcommand{\bx}{\mathbf{x}}
\newcommand{\by}{\mathbf{y}}
\newcommand{\ba}{\mathbf{a}}

\title{From Least Action to Diffusion and Hydrodynamics:\\
A Unified Variational Perspective}
\author{}
\date{\today}

\begin{document}
\maketitle

\begin{abstract}
\noindent
We show how two paradigmatic equations of continuum physics -- the diffusion equation and the Euler equations of hydrodynamics -- can both be obtained from extremal, ``least-action''-type principles, despite the fact that one is dissipative and irreversible and the other is not. For an ideal barotropic fluid, Hamilton's principle applies directly once the Eulerian velocity is represented through a velocity potential, and free variation of a single action reproduces both the continuity equation and Euler's equation. For the diffusion equation we first show, via a self-adjointness (Helmholtz) argument, why no stationary-action principle can exist in terms of the density alone; we then exhibit two constructions that do produce it: a doubled-field action in the manner of Bateman and Morse--Feshbach, and a genuine minimization principle -- the Wasserstein-metric gradient flow of the Boltzmann entropy, in the sense of Jordan, Kinderlehrer and Otto. We close by showing that both constructions are special cases of extremizing one family of energy functionals (kinetic energy plus a Fisher-information term) on the space of probability densities: a Hamiltonian (symplectic) extremization reproduces the quantum-hydrodynamic (Madelung) equations, while a metric-gradient extremization of the closely related entropy functional reproduces the diffusion equation. The two families of equations are thus not merely analogous, but are the reversible and irreversible flows generated by one underlying variational structure.
\end{abstract}

\section{Introduction}

Hamilton's principle of stationary action, $\delta\int L\,\rd t = 0$, is usually first met in the setting of finite-dimensional, time-reversible mechanical systems. Two extensions of this idea are less commonly spelled out side by side. First, continuum systems with infinitely many degrees of freedom admit an entirely analogous field-theoretic formulation, provided the Eulerian description is handled with some care. Second, and more surprisingly, genuinely irreversible equations can also be brought under an extremal-principle umbrella -- but only by either enlarging the space of fields (doubling), or by replacing \emph{stationarity} with \emph{minimization} with respect to a metric structure that is not symplectic.

This paper works out both extensions for two concrete equations. Section~\ref{sec:hydro} derives the continuity and Euler equations of an ideal, barotropic fluid from Hamilton's principle, following the classical program of Bateman, Clebsch, Eckart, Luke, and Seliger and Whitham \cite{Bateman1929,Clebsch1859,Eckart1960,Luke1967,SeligerWhitham1968}. Section~\ref{sec:diffusion} turns to the diffusion equation $\partial_t\rho = D\lap\rho$: we show precisely why it \emph{cannot} arise from a bare stationary-action principle (an instance of the classical inverse problem of the calculus of variations \cite{Bauer1931,Bateman1931,MorseFeshbach1953}), and then derive it from two different extremal constructions that evade the obstruction. Section~\ref{sec:unification} shows that hydrodynamics and diffusion are not simply analogous constructions but arise as the Hamiltonian and gradient flows, respectively, of one family of functionals on the space of probability densities -- the bridge being the Madelung transformation of the Schr\"odinger equation \cite{Madelung1927} and Nelson's stochastic mechanics \cite{Nelson1966,GuerraMorato1983,Reginatto1998}, made geometrically precise by optimal-transport (Wasserstein) theory \cite{JordanKinderlehrerOtto1998,Otto2001,BenamouBrenier2000,vonRenesse2012,GianazzaSavareToscani2009}. Section~\ref{sec:conclusion} summarizes.

\section{Hydrodynamic equations from Hamilton's principle}
\label{sec:hydro}

\subsection{The Eulerian obstruction}

Consider an ideal, compressible, barotropic fluid. In the \emph{material} (Lagrangian) description, a fluid parcel initially at $\ba$ occupies position $\bx(\ba,t)$ at time $t$, and the action is simply that of a continuum of free particles moving in a potential built from an internal-energy density:
\begin{equation}
S_{\mathrm{mat}}[\bx] = \int \rd t \int \rd^3 a\; \rho_0(\ba)\left[\frac12\left|\pd{\bx}{t}\right|^2 - \Phi(\bx(\ba,t))\right] - \int \rd t\, U[\rho],
\end{equation}
where $\rho_0$ is the initial density and $U[\rho]=\int w(\rho)\,\rd^3x$ is the total internal energy, $\rho$ being determined kinematically by mass conservation, $\rho(\bx(\ba,t),t)\,\det(\partial x_i/\partial a_j) = \rho_0(\ba)$. Free variation of $\bx(\ba,t)$ here reproduces Newton's second law parcel by parcel and is unproblematic.

The difficulty, pointed out already by Bateman \cite{Bateman1929} and analyzed systematically by Seliger and Whitham \cite{SeligerWhitham1968}, is that one would like a variational principle stated directly in terms of the \emph{Eulerian} fields $\rho(\bx,t)$ and $\bv(\bx,t)$. A naive attempt to vary $\rho$ and $\bv$ independently fails, because at fixed $\bx$ these are not independent: a displacement of fluid parcels changes both, in a way that a bare Eulerian variation does not see. The resolution is to introduce potentials for the velocity field so that the variation is again taken with respect to genuinely independent functions.

\subsection{A working Eulerian action}

For potential (irrotational) barotropic flow, write $\bv = \nabla\phi$ and treat $\rho(\bx,t)$ and $\phi(\bx,t)$ as independent fields. Let $w(\rho)$ denote the internal energy per unit volume, and define the pressure by the standard barotropic relation
\begin{equation}
p(\rho) \equiv \rho\, w'(\rho) - w(\rho). \label{eq:pressure-def}
\end{equation}
Consider the action (essentially due to Bateman \cite{Bateman1929} and Luke \cite{Luke1967}; see also Eckart \cite{Eckart1960})
\begin{equation}
S[\rho,\phi] = \int \rd t \int \rd^3x \; \mathcal{L}, \qquad
\mathcal{L} = -\rho\,\partial_t\phi - \frac12\rho\,|\nabla\phi|^2 - \rho\,\Phi(\bx) - w(\rho). \label{eq:hydro-action}
\end{equation}
This is not simply ``kinetic minus potential energy'': the first term plays the role of a $p\,\dot q$ term, making $\rho$ and $\phi$ a canonically conjugate pair, exactly as in the Hamiltonian (Clebsch/Zakharov) formulation of fluid mechanics.

\paragraph{Variation with respect to $\phi$.} Since $\phi$ enters $\mathcal L$ only through $\partial_t\phi$ and $\nabla\phi$, the Euler--Lagrange equation is
\begin{equation}
-\partial_t\left(\frac{\partial\mathcal L}{\partial(\partial_t\phi)}\right) - \nabla\cdot\left(\frac{\partial \mathcal L}{\partial(\nabla\phi)}\right) = 0
\;\Longrightarrow\;
\partial_t\rho + \nabla\cdot(\rho\nabla\phi) = 0,
\end{equation}
which is precisely the \textbf{continuity equation} $\partial_t\rho + \nabla\cdot(\rho\bv)=0$ with $\bv=\nabla\phi$.

\paragraph{Variation with respect to $\rho$.} Here $\mathcal L$ depends on $\rho$ only algebraically, so
\begin{equation}
\frac{\partial \mathcal L}{\partial \rho} = -\partial_t\phi - \frac12|\nabla\phi|^2 - \Phi - w'(\rho) = 0
\;\Longrightarrow\;
\partial_t \phi + \frac12|\nabla\phi|^2 + \Phi + w'(\rho) = 0, \label{eq:bernoulli}
\end{equation}
the unsteady \textbf{Bernoulli equation}. Taking $\nabla$ of \eqref{eq:bernoulli} and using $\bv=\nabla\phi$,
\begin{equation}
\partial_t \bv + \nabla\!\left(\tfrac12|\bv|^2\right) + \nabla\Phi + w''(\rho)\nabla\rho = 0.
\end{equation}
Because the flow is irrotational, $(\bv\cdot\nabla)\bv = \nabla(\tfrac12|\bv|^2) - \bv\times(\nabla\times\bv) = \nabla(\tfrac12|\bv|^2)$; and from the definition \eqref{eq:pressure-def}, $\nabla p = \rho\,w''(\rho)\nabla\rho$, so $w''(\rho)\nabla\rho = \nabla p/\rho$. Substituting,
\begin{equation}
\boxed{\;\rho\left[\partial_t \bv + (\bv\cdot\nabla)\bv\right] = -\nabla p - \rho\nabla\Phi\;} \label{eq:euler}
\end{equation}
which is precisely \textbf{Euler's equation} for a barotropic, inviscid fluid. Equations \eqref{eq:bernoulli}--\eqref{eq:euler} together with the continuity equation constitute the full hydrodynamic system, obtained from the single stationary-action principle $\delta S[\rho,\phi]=0$.

\subsection{Remark: vorticity and Clebsch potentials}

The representation $\bv=\nabla\phi$ excludes vorticity. The general case is recovered by the Clebsch representation $\bv = \nabla\phi + \alpha\nabla\beta$, with an extra term $-\rho\alpha\,\partial_t\beta$ added to $\mathcal{L}$ in \eqref{eq:hydro-action}; free variation with respect to the four independent potentials $(\phi,\rho,\alpha,\beta)$ then reproduces the general (rotational) Euler equation, with $\alpha,\beta$ playing the role of a pair of scalars advected by the flow, $D\alpha/Dt = D\beta/Dt = 0$. This construction, due to Clebsch \cite{Clebsch1859} and put on a systematic footing by Seliger and Whitham \cite{SeligerWhitham1968}, is not needed below and we do not pursue it further; we record it only to emphasize that the potential-flow case treated above is the simplest instance of a more general Eulerian variational principle, not a special trick.

\section{The diffusion equation and the limits of Hamilton's principle}
\label{sec:diffusion}

\subsection{An obstruction: self-adjointness}
\label{sec:obstruction}

It is natural to ask whether the diffusion equation
\begin{equation}
\partial_t \rho = D\lap \rho \label{eq:diffusion}
\end{equation}
can likewise be obtained as $\delta S[\rho]/\delta\rho = 0$ for some action functional of $\rho$ alone. The classical (Helmholtz) inverse problem of the calculus of variations answers this in the negative, and the reason is instructive.

A necessary condition for a system of equations $\mathcal N[\rho]=0$ to arise as the Euler--Lagrange equation of \emph{some} action $S[\rho]$ is that its linearization (here, since \eqref{eq:diffusion} is already linear, the operator itself) be formally self-adjoint with respect to the natural pairing $\langle f,g\rangle = \int_{t_0}^{t_1}\!\!\int f g\;\rd^3x\,\rd t$. Write $\mathcal O \equiv \partial_t - D\lap$. For compactly supported (or suitably decaying) test functions $f,g$, two integrations by parts in space and one in time give
\begin{equation}
\langle f, \mathcal O g\rangle = \int f(\partial_t g - D\lap g) = \int (-\partial_t f - D\lap f)\,g = \langle \mathcal O^{\dagger} f, g\rangle, \qquad \mathcal O^{\dagger} = -\partial_t - D\lap.
\end{equation}
Since $\mathcal O^\dagger \neq \mathcal O$ (the spatial part is self-adjoint but the time derivative changes sign under integration by parts while $D\lap$ does not), $\mathcal O$ is \emph{not} self-adjoint, and no action $S[\rho]$ exists whose stationarity gives \eqref{eq:diffusion}. This is a special case of a general theorem of Bauer \cite{Bauer1931}: a linear equation with constant coefficients that is not invariant under $t\to -t$ cannot be the Euler--Lagrange equation of an action depending on the field alone. Physically, this is just the statement that stationary-action dynamics is always time-reversible in the appropriate sense, while diffusion manifestly is not.

Bateman \cite{Bateman1931} showed how to evade this obstruction in general: enlarge the field content by introducing a second, ``mirror'' field that absorbs exactly what the first field dissipates. We apply this to \eqref{eq:diffusion} next; see also Morse and Feshbach \cite{MorseFeshbach1953}.

\subsection{A doubled-field least action principle}
\label{sec:doubled}

Introduce an auxiliary field $\tilde\rho(\bx,t)$ and the action
\begin{equation}
S[\rho,\tilde\rho] = \int \rd t\int \rd^3x\; \tilde\rho\left(\partial_t \rho - D\lap\rho\right). \label{eq:doubled-action}
\end{equation}

\paragraph{Variation with respect to $\tilde\rho$.} Since \eqref{eq:doubled-action} is linear in $\tilde\rho$, the Euler--Lagrange equation is simply
\begin{equation}
\partial_t \rho - D\lap \rho = 0,
\end{equation}
the diffusion equation itself.

\paragraph{Variation with respect to $\rho$.} Integrating the first term by parts in $t$ and the second twice by parts in space (boundary terms discarded),
\begin{equation}
\int \tilde\rho\,\partial_t\rho = -\int (\partial_t\tilde\rho)\,\rho, \qquad \int \tilde\rho\,(-D\lap\rho) = \int(-D\lap\tilde\rho)\,\rho,
\end{equation}
so the Euler--Lagrange equation with respect to $\rho$ is
\begin{equation}
-\partial_t\tilde\rho - D\lap\tilde\rho = 0 \quad\Longleftrightarrow\quad \partial_t\tilde\rho = -D\lap\tilde\rho. \label{eq:adjoint-diffusion}
\end{equation}
This is exactly the \emph{backward} (time-reversed, or ``anti-diffusion'') equation: setting $\tau=-t$, \eqref{eq:adjoint-diffusion} becomes an ordinary forward diffusion equation for $\tilde\rho$ in $\tau$. The pair $(\rho,\tilde\rho)$ is precisely of Bateman's type \cite{Bateman1931}: $\rho$ dissipates (spreads and loses gradients) while $\tilde\rho$ is the complementary field that would absorb exactly what $\rho$ loses, run in reverse time. In the language later used for stochastic dynamics (Onsager--Machlup, Martin--Siggia--Rose), $\tilde\rho$ plays the role of a \emph{response field}: it has no direct physical density interpretation but its presence is unavoidable if one insists on a strict $\delta S=0$ formulation. This doubling is not a peculiarity of the diffusion equation; it is the same device used for the textbook example of a damped harmonic oscillator, which likewise has no one-field Lagrangian \cite{Bateman1931}.

\subsection{A genuine minimization: gradient flow in Wasserstein space}
\label{sec:jko}

The doubled-field construction produces a \emph{stationary}-action principle, at the cost of a field with no independent physical meaning. It is natural to ask whether \eqref{eq:diffusion} is instead the outcome of an honest \emph{minimization} -- the sense in which ``least action'' is most literally read. Onsager's variational principle for near-equilibrium irreversible processes \cite{Onsager1931I,Onsager1931II} already suggested that dissipative evolution extremizes a dissipation-type functional rather than Hamilton's action; the sharp, modern form of this idea for the diffusion equation is due to Jordan, Kinderlehrer and Otto \cite{JordanKinderlehrerOtto1998}, with the general geometric framework supplied by Otto \cite{Otto2001}.

Let $\mathcal P(\mathbb R^d)$ denote probability densities on $\mathbb R^d$ and let
\begin{equation}
W_2(\mu,\nu)^2 = \min_{\pi\in\Pi(\mu,\nu)} \int |\bx-\by|^2\,\rd\pi(\bx,\by)
\end{equation}
be the quadratic Wasserstein (optimal-transport) distance, the minimum being over couplings $\pi$ with marginals $\mu,\nu$. Define the Boltzmann entropy functional
\begin{equation}
\mathcal F[\rho] = D\int \rho\log\rho\;\rd^3x.
\end{equation}
The Jordan--Kinderlehrer--Otto (JKO) \emph{minimizing-movement} scheme discretizes time into steps of size $\tau$ and defines $\rho^{n+1}$ as the literal minimizer
\begin{equation}
\rho^{n+1} = \operatorname*{arg\,min}_{\rho\in\mathcal P(\mathbb R^d)} \left\{\frac{1}{2\tau}\,W_2(\rho,\rho^n)^2 + \mathcal F[\rho]\right\}. \label{eq:jko}
\end{equation}
This is a genuine least-action principle at each step: it minimizes a weighted sum of a ``transport cost'' from the previous state and the entropy of the new one. Jordan, Kinderlehrer and Otto proved \cite{JordanKinderlehrerOtto1998} that as $\tau\to 0$ the piecewise-constant interpolation of $\{\rho^n\}$ converges to the solution of the Wasserstein-metric gradient flow of $\mathcal F$,
\begin{equation}
\partial_t \rho = \nabla\cdot\left(\rho\,\nabla\frac{\delta \mathcal F}{\delta\rho}\right). \label{eq:wgf}
\end{equation}
Since $\delta\mathcal F/\delta\rho = D(\log\rho + 1)$, we have $\nabla(\delta\mathcal F/\delta\rho) = D\,\nabla\rho/\rho$, and \eqref{eq:wgf} becomes
\begin{equation}
\partial_t\rho = \nabla\cdot\left(\rho\cdot D\frac{\nabla\rho}{\rho}\right) = D\lap\rho,
\end{equation}
recovering the diffusion equation \eqref{eq:diffusion} exactly -- this time as the genuine minimization \eqref{eq:jko}, with no auxiliary field required, at the price of using a metric (Riemannian, non-symplectic) rather than symplectic extremal structure on the space of densities.

It is worth noting the family resemblance to Section~\ref{sec:hydro}. Benamou and Brenier \cite{BenamouBrenier2000} showed that $W_2$ itself has a dynamic, action-like representation,
\begin{equation}
W_2(\rho_0,\rho_1)^2 = \min_{\rho,\bv} \left\{ \int_0^1\!\!\int \rho\,|\bv|^2\,\rd^3x\,\rd t \;:\; \partial_t\rho + \nabla\cdot(\rho\bv)=0,\ \rho(0)=\rho_0,\ \rho(1)=\rho_1\right\},
\end{equation}
i.e., precisely a least kinetic-action problem constrained by the continuity equation, structurally the same object as the kinetic term in \eqref{eq:hydro-action}. The JKO scheme \eqref{eq:jko} can be read as this same kinetic action, with the entropy $\mathcal F$ inserted step-by-step as a relaxation term. Sections~\ref{sec:hydro} and \ref{sec:jko} are thus two different extremizations -- one genuinely Hamiltonian, one a metric gradient flow -- built from the same underlying kinetic functional. We make this precise in Section~\ref{sec:unification}.

\section{Unifying perspective: Hamiltonian versus gradient flow on the space of densities}
\label{sec:unification}

Return to the Eulerian action \eqref{eq:hydro-action} and add to it the term $-\dfrac{\hbar^2}{2m}(\nabla\sqrt\rho)^2$, writing $S$ for the potential $m\phi$ so as to match standard quantum-mechanical conventions:
\begin{equation}
\mathcal L_Q = -\rho\,\partial_t S - \frac{\rho}{2m}|\nabla S|^2 - \rho V(\bx) - \frac{\hbar^2}{2m}(\nabla\sqrt\rho)^2. \label{eq:madelung-action}
\end{equation}
Variation with respect to $S$ again gives the continuity equation, $\partial_t\rho + \nabla\cdot(\rho\nabla S/m)=0$. Variation with respect to $\rho$ -- using $(\nabla\sqrt\rho)^2 = |\nabla\rho|^2/4\rho$ and the identity $\dfrac{\hbar^2}{2m}\dfrac{\lap\sqrt\rho}{\sqrt\rho} = \dfrac{\hbar^2}{4m}\dfrac{\lap\rho}{\rho} - \dfrac{\hbar^2}{8m}\dfrac{|\nabla\rho|^2}{\rho^2}$, which one verifies directly by differentiating $\sqrt\rho$ twice -- gives
\begin{equation}
\partial_t S + \frac{|\nabla S|^2}{2m} + V + Q(\rho) = 0, \qquad Q(\rho) \equiv -\frac{\hbar^2}{2m}\frac{\lap\sqrt\rho}{\sqrt\rho}. \label{eq:madelung-hj}
\end{equation}
Equations for $(\rho,S)$ above are exactly the \textbf{Madelung equations} \cite{Madelung1927}: setting $\psi = \sqrt\rho\,e^{iS/\hbar}$, they are equivalent to the free Schr\"odinger equation $i\hbar\,\partial_t\psi = -\dfrac{\hbar^2}{2m}\lap\psi + V\psi$, with $Q(\rho)$ the Bohm quantum potential. So \eqref{eq:madelung-action} is simply the hydrodynamic action \eqref{eq:hydro-action} with the fluid's internal energy $w(\rho)$ replaced by $\dfrac{\hbar^2}{2m}(\nabla\sqrt\rho)^2$ -- a quantum ``pressure'' built not from $\rho$ itself but from its gradient.

The quantity $I[\rho] \equiv \displaystyle\int \frac{|\nabla\rho|^2}{\rho}\,\rd^3x$ is the (classical) \emph{Fisher information} of $\rho$, and $\dfrac{\hbar^2}{2m}(\nabla\sqrt\rho)^2 = \dfrac{\hbar^2}{8m}\dfrac{|\nabla\rho|^2}{\rho}$, so the extra term in \eqref{eq:madelung-action} is $-\dfrac{\hbar^2}{8m}$ times the Fisher-information density. That the Schr\"odinger equation can be obtained by extremizing kinetic energy plus potential energy plus a fixed multiple of the Fisher information is not a coincidence internal to this derivation: it is also the content of Reginatto's principle of minimum Fisher information \cite{Reginatto1998} and of Hall and Reginatto's exact-uncertainty derivation, and it is made fully rigorous, in the language of formal Riemannian (Otto) calculus on the Wasserstein space $\mathcal P_2(\mathbb R^d)$, by von Renesse \cite{vonRenesse2012}: the Schr\"odinger equation is exactly Newton's second law, $\nabla^{\mathcal W}_{\dot\mu}\dot\mu = -\nabla^{\mathcal W}F(\mu)$, on $\mathcal P_2(\mathbb R^d)$ equipped with the Wasserstein Riemannian metric, for the potential $F(\mu) = \langle V,\mu\rangle + \frac{\hbar^2}{8}\int|\nabla\log\mu|^2\,\rd\mu$ -- i.e., a genuinely \emph{Hamiltonian} (symplectic) flow generated by exactly this kinetic-plus-Fisher-information energy.

Now compare this with Section~\ref{sec:jko}. There, the diffusion equation was obtained as the \emph{gradient} (metric, non-symplectic) flow of the entropy $\mathcal F[\rho]=D\int\rho\log\rho$ on the very same space $\mathcal P_2(\mathbb R^d)$. Gianazza, Savar\'e and Toscani \cite{GianazzaSavareToscani2009} close the loop directly: they show that the Fisher information functional itself, $\mathcal F^f[\rho] = \tfrac12\int|\nabla\log\rho|^2\rho\,\rd^3x + \int f\rho\,\rd^3x$, generates -- again as a \emph{Wasserstein gradient flow}, in the sense of Section~\ref{sec:jko} -- a fourth-order parabolic equation, the \emph{quantum drift-diffusion equation}, $\partial_t \rho + \nabla\cdot\!\left(\rho\, D\!\left(2\dfrac{\lap\sqrt\rho}{\sqrt\rho} - f\right)\right) = 0$: precisely the diffusive (dissipative, metric-gradient) counterpart of the Hamiltonian quantum-hydrodynamic system \eqref{eq:madelung-hj}, built from the same Bohm/Fisher term $\lap\sqrt\rho/\sqrt\rho$.

The overall picture is therefore:
\begin{itemize}
\item[--] The \emph{same} kinetic functional $\int \rho|\bv|^2$, extremized subject to the continuity constraint, is Hamilton's principle for pressureless potential flow (Section~\ref{sec:hydro}) and, in its dynamic (Benamou--Brenier) form, generates the Wasserstein metric on $\mathcal P_2(\mathbb R^d)$ (Section~\ref{sec:jko}).
\item[--] A \emph{Hamiltonian} (symplectic) extremization of kinetic $+$ potential $+$ Fisher-information energy on $\mathcal P_2(\mathbb R^d)$ gives the Madelung/Schr\"odinger system \eqref{eq:madelung-hj} (this section).
\item[--] A \emph{metric-gradient} extremization of the closely related entropy, or of the Fisher information itself, on the same space gives, respectively, ordinary diffusion (Section~\ref{sec:jko}) and the quantum drift-diffusion equation \cite{GianazzaSavareToscani2009}.
\end{itemize}
Hydrodynamics/quantum mechanics and diffusion are, in this precise sense, the reversible and irreversible flows of one underlying variational structure on the space of probability densities, rather than two unrelated uses of the phrase ``least action.''

As a final, more heuristic remark: Nelson's stochastic mechanics \cite{Nelson1966} and the Guerra--Morato stochastic variational principle \cite{GuerraMorato1983} arrive at the same correspondence from the probabilistic side, by exhibiting the Schr\"odinger equation as the outcome of a stochastic least-action principle for a diffusion process with diffusion coefficient $D=\hbar/2m$. Formally continuing $D \to i\hbar/2m$ (equivalently, Wick-rotating $t\to -i\tau$) turns the diffusion equation $\partial_\tau\rho = D\lap\rho$ into the free Schr\"odinger equation. This substitution should be treated with some caution: it exchanges a real, positivity-preserving parabolic equation for a complex, unitary, dispersive one, and the well-posedness theory of the two is quite different. It is nonetheless the precise sense, familiar from Euclidean quantum mechanics, in which diffusion and quantum hydrodynamics are analytic continuations of each other -- consistent with, and a compact restatement of, the variational picture developed above.

\section{Conclusion}
\label{sec:conclusion}

We have derived the hydrodynamic equations of an ideal barotropic fluid, and the diffusion equation, from extremal principles, while being explicit about where the two derivations differ in kind. (i) The Euler and continuity equations follow from a bona fide Hamilton's stationary-action principle once the Eulerian velocity is represented through a velocity potential (or, more generally, Clebsch potentials); no extra structure is needed. (ii) The diffusion equation, being genuinely dissipative, cannot arise from a stationary-action principle in the density alone -- a consequence of the non-self-adjointness of the diffusion operator -- but it does follow from two related constructions: a doubled-field stationary action in the manner of Bateman, and a literal minimization, the Wasserstein-metric gradient flow of the Boltzmann entropy. (iii) Both hydrodynamics (in its quantum, Madelung form) and diffusion turn out to be flows of the same family of kinetic-plus-Fisher-information functionals on the space of probability densities -- a Hamiltonian flow in the first case, a metric gradient flow in the second -- with the Schr\"odinger equation and Nelson's stochastic mechanics providing the historical and conceptual bridge between them.

\begin{thebibliography}{99}

\bibitem{Bateman1929} H.~Bateman, ``Notes on a differential equation which occurs in the two-dimensional motion of a compressible fluid and the associated variational problems,'' Proc.\ R.\ Soc.\ Lond.\ A \textbf{125}, 598--618 (1929).

\bibitem{Clebsch1859} A.~Clebsch, ``Ueber die Integration der hydrodynamischen Gleichungen,'' J.\ Reine Angew.\ Math.\ \textbf{56}, 1--10 (1859).

\bibitem{Eckart1960} C.~Eckart, ``Variation Principles of Hydrodynamics,'' Phys.\ Fluids \textbf{3}, 421--427 (1960).

\bibitem{Luke1967} J.~C.~Luke, ``A variational principle for a fluid with a free surface,'' J.\ Fluid Mech.\ \textbf{27}, 395--397 (1967).

\bibitem{SeligerWhitham1968} R.~L.~Seliger and G.~B.~Whitham, ``Variational Principles in Continuum Mechanics,'' Proc.\ R.\ Soc.\ Lond.\ A \textbf{305}, 1--25 (1968).

\bibitem{Bauer1931} P.~S.~Bauer, ``Dissipative Dynamical Systems I,'' Proc.\ Natl.\ Acad.\ Sci.\ USA \textbf{17}, 311--314 (1931).

\bibitem{Bateman1931} H.~Bateman, ``On Dissipative Systems and Related Variational Principles,'' Phys.\ Rev.\ \textbf{38}, 815--819 (1931).

\bibitem{MorseFeshbach1953} P.~M.~Morse and H.~Feshbach, \emph{Methods of Theoretical Physics} (McGraw-Hill, New York, 1953).

\bibitem{Onsager1931I} L.~Onsager, ``Reciprocal Relations in Irreversible Processes. I,'' Phys.\ Rev.\ \textbf{37}, 405--426 (1931).

\bibitem{Onsager1931II} L.~Onsager, ``Reciprocal Relations in Irreversible Processes. II,'' Phys.\ Rev.\ \textbf{38}, 2265--2279 (1931).

\bibitem{JordanKinderlehrerOtto1998} R.~Jordan, D.~Kinderlehrer, and F.~Otto, ``The Variational Formulation of the Fokker--Planck Equation,'' SIAM J.\ Math.\ Anal.\ \textbf{29}(1), 1--17 (1998).

\bibitem{Otto2001} F.~Otto, ``The geometry of dissipative evolution equations: the porous medium equation,'' Comm.\ Partial Differential Equations \textbf{26}(1--2), 101--174 (2001).

\bibitem{BenamouBrenier2000} J.-D.~Benamou and Y.~Brenier, ``A computational fluid mechanics solution to the Monge--Kantorovich mass transfer problem,'' Numer.\ Math.\ \textbf{84}(3), 375--393 (2000).

\bibitem{Madelung1927} E.~Madelung, ``Quantentheorie in hydrodynamischer Form,'' Z.\ Phys.\ \textbf{40}, 322--326 (1927).

\bibitem{Nelson1966} E.~Nelson, ``Derivation of the Schr\"odinger Equation from Newtonian Mechanics,'' Phys.\ Rev.\ \textbf{150}, 1079--1085 (1966).

\bibitem{GuerraMorato1983} F.~Guerra and L.~M.~Morato, ``Quantization of dynamical systems and stochastic control theory,'' Phys.\ Rev.\ D \textbf{27}(8), 1774--1786 (1983).

\bibitem{Reginatto1998} M.~Reginatto, ``Derivation of the equations of nonrelativistic quantum mechanics using the principle of minimum Fisher information,'' Phys.\ Rev.\ A \textbf{58}, 1775--1778 (1998).

\bibitem{vonRenesse2012} M.-K.~von Renesse, ``An Optimal Transport View of Schr\"odinger's Equation,'' Canad.\ Math.\ Bull.\ \textbf{55}(4), 858--869 (2012).

\bibitem{GianazzaSavareToscani2009} U.~Gianazza, G.~Savar\'e, and G.~Toscani, ``The Wasserstein Gradient Flow of the Fisher Information and the Quantum Drift-Diffusion Equation,'' Arch.\ Ration.\ Mech.\ Anal.\ \textbf{194}(1), 133--220 (2009).

\end{thebibliography}

\end{document}
